% -*- mode : latex; coding: utf-8 -*-
\markboth{Abstract}{}
\begin{EnglishAbstract}
  RowHammer is a serious security threat to computing systems using DRAM since the first 
  experimental analysis in 2014 \cite{FlippingBits}. Through repetitive activation of the DRAM row, 
  electromagnetic inference causes charge leakage in DRAM cells. Even worse, technology scaling 
  makes DRAM even more vulnerable to RowHammer attack \cite{Revisiting}. 
  Many RowHammer mitigation systems have been proposed, but a lack of understanding of the internal DRAM structure 
  poses inefficiency in the system.
  In this paper, I point out some optimization points to enhance the 
  current state-of-the-art RowHammer mitigation schemes using internal DRAM topologies in DDR4 memory system: 
  Graphene \cite{Graphene} and Hydra \cite{Hydra}. I evaluate speedup between baseline and optimized 
  version in Macsim with Ramulator \cite{MACSIM}, \cite{RAMULATOR} (cycle-accurate and trace-driven 
  simulator, with 2-channel DDR4 configuration) with 2 synthetic adversarial attack patterns and 3 
  memory-intensive workloads from SPEC CPU 2017 under the RowHammer threshold of 1K. 
  Results show that with 4\% additional CAM in Graphene, an additional refresh was reduced by 
  nearly 50\% in adversarial patterns and 64.7\% in normal workloads with a performance overhead reduced by 65.7\%. 
  In the case of Hydra, an additional refresh was reduced by 91.4\% in adversarial pattern and 98.1\% in normal workloads 
  with average performance overhead reduced by 71.7\%. 
\end{EnglishAbstract}
